\documentclass[10pt,a4paper]{amsart}
\usepackage{amsmath,amssymb,amsthm,graphicx}
% Retained minimal preamble for the Stacks Project chapter varieties.tex.
% Verbatim excerpt of the chapter preamble: Xy-pic setup (native preamble lines 14-18), the native theorem
% declarations with the mathematical macros (native preamble lines 154-217), and the project own no-op
% definitions of the reference/slogan/history display environments (native wrapper.tex lines 1-3) that the
% counted statement uses for its slogan. Copied mechanically; nothing added, removed or redefined.
% For commutative diagrams we use Xy-pic
\usepackage[all]{xy}

% We use 2cell for 2-commutative diagrams.
\xyoption{2cell}
% Theorem environments.
%
\theoremstyle{plain}
\newtheorem{theorem}[subsection]{Theorem}
\newtheorem{proposition}[subsection]{Proposition}
\newtheorem{lemma}[subsection]{Lemma}

\theoremstyle{definition}
\newtheorem{definition}[subsection]{Definition}
\newtheorem{example}[subsection]{Example}
\newtheorem{exercise}[subsection]{Exercise}
\newtheorem{situation}[subsection]{Situation}

\theoremstyle{remark}
\newtheorem{remark}[subsection]{Remark}
\newtheorem{remarks}[subsection]{Remarks}

\numberwithin{equation}{subsection}

% Macros
%
\def\lim{\mathop{\mathrm{lim}}\nolimits}
\def\colim{\mathop{\mathrm{colim}}\nolimits}
\def\Spec{\mathop{\mathrm{Spec}}}
\def\Hom{\mathop{\mathrm{Hom}}\nolimits}
\def\Ext{\mathop{\mathrm{Ext}}\nolimits}
\def\SheafHom{\mathop{\mathcal{H}\!\mathit{om}}\nolimits}
\def\SheafExt{\mathop{\mathcal{E}\!\mathit{xt}}\nolimits}
\def\Sch{\mathit{Sch}}
\def\Mor{\mathop{\mathrm{Mor}}\nolimits}
\def\Ob{\mathop{\mathrm{Ob}}\nolimits}
\def\Sh{\mathop{\mathit{Sh}}\nolimits}
\def\NL{\mathop{N\!L}\nolimits}
\def\CH{\mathop{\mathrm{CH}}\nolimits}
\def\proetale{{pro\text{-}\acute{e}tale}}
\def\etale{{\acute{e}tale}}
\def\QCoh{\mathit{QCoh}}
\def\Ker{\mathop{\mathrm{Ker}}}
\def\Im{\mathop{\mathrm{Im}}}
\def\Coker{\mathop{\mathrm{Coker}}}
\def\Coim{\mathop{\mathrm{Coim}}}

% Boxtimes
%
\DeclareMathSymbol{\boxtimes}{\mathbin}{AMSa}{"02}

%
% Macros for moduli stacks/spaces
%
\def\QCohstack{\mathcal{QC}\!\mathit{oh}}
\def\Cohstack{\mathcal{C}\!\mathit{oh}}
\def\Spacesstack{\mathcal{S}\!\mathit{paces}}
\def\Quotfunctor{\mathrm{Quot}}
\def\Hilbfunctor{\mathrm{Hilb}}
\def\Curvesstack{\mathcal{C}\!\mathit{urves}}
\def\Polarizedstack{\mathcal{P}\!\mathit{olarized}}
\def\Complexesstack{\mathcal{C}\!\mathit{omplexes}}
% \Pic is the operator that assigns to X its picard group, usage \Pic(X)
% \Picardstack_{X/B} denotes the Picard stack of X over B
% \Picardfunctor_{X/B} denotes the Picard functor of X over B
\def\Pic{\mathop{\mathrm{Pic}}\nolimits}
\def\Picardstack{\mathcal{P}\!\mathit{ic}}
\def\Picardfunctor{\mathrm{Pic}}
\def\Deformationcategory{\mathcal{D}\!\mathit{ef}}
\newenvironment{reference}{}{}
\newenvironment{slogan}{}{}
\newenvironment{history}{}{}
\usepackage{cleveref}
\crefname{originallemma}{original lemma}{original lemmas}
\Crefname{originallemma}{Original lemma}{Original lemmas}
\makeatletter
\newlabel{lemma-normalization-and-change-of-fields}{{0C3N}{source}{}{sourcepointer-1}{}}
\newlabel{lemma-normalization-and-change-of-fields@cref}{{[originallemma][0][]0C3N}{[1][0][]source}{}{}{}}
\makeatother
\crefname{originalreference}{original reference}{original references}
\Crefname{originalreference}{Original reference}{Original references}
\makeatletter
\newlabel{morphisms-lemma-universal-homeomorphism}{{04DF}{source}{}{sourcepointer-2}{}}
\newlabel{morphisms-lemma-universal-homeomorphism@cref}{{[originalreference][0][]04DF}{[1][0][]source}{}{}{}}
\makeatother
\crefname{originalreference}{original reference}{original references}
\Crefname{originalreference}{Original reference}{Original references}
\makeatletter
\newlabel{morphisms-lemma-universally-injective}{{01S4}{source}{}{sourcepointer-3}{}}
\newlabel{morphisms-lemma-universally-injective@cref}{{[originalreference][0][]01S4}{[1][0][]source}{}{}{}}
\makeatother
\crefname{originallemma}{original lemma}{original lemmas}
\Crefname{originallemma}{Original lemma}{Original lemmas}
\makeatletter
\newlabel{lemma-number-of-branches}{{0C1S}{source}{}{sourcepointer-4}{}}
\newlabel{lemma-number-of-branches@cref}{{[originallemma][0][]0C1S}{[1][0][]source}{}{}{}}
\makeatother
\crefname{originalreference}{original reference}{original references}
\Crefname{originalreference}{Original reference}{Original references}
\makeatletter
\newlabel{schemes-lemma-points-fibre-product}{{01JT}{source}{}{sourcepointer-5}{}}
\newlabel{schemes-lemma-points-fibre-product@cref}{{[originalreference][0][]01JT}{[1][0][]source}{}{}{}}
\makeatother
\crefname{originallemma}{original lemma}{original lemmas}
\Crefname{originallemma}{Original lemma}{Original lemmas}
\makeatletter
\newlabel{lemma-separably-closed-field-connected-components}{{0363}{source}{}{sourcepointer-6}{}}
\newlabel{lemma-separably-closed-field-connected-components@cref}{{[originallemma][0][]0363}{[1][0][]source}{}{}{}}
\makeatother
\title{The number of geometric branches at a point is invariant under extension of the base field}
\author{The Stacks Project Authors}
\date{stacks-project:varieties.tex:local-snapshot-20261011; GNU FDL 1.2 or later}
\begin{document}
\maketitle
\noindent\textit{Editorial note.} All delivered mathematical text is the author's own, copied verbatim from the pinned Stacks Project chapter file varieties.tex; no statement or proof step was rewritten, reordered or removed.\par\medskip
\noindent\textit{Editorial note.} The proof contains a commutative diagram typeset with Xy-pic; the retained chapter preamble supplies the native Xy-pic setup, so the diagram compiles exactly as in the source.\par\medskip
\noindent\textit{Editorial note.} The counted span cites six cross-chapter labels, each recorded in assembly.external\_references with its canonical tag as a navigation pointer.\par\medskip
\noindent\textit{Editorial note.} The final algebra fact is attributed by the author to an elementary argument as well as to the cited lemma; the counted span contains the author's own reduction to that fact.\par\medskip
\noindent\textit{Allowed standard prerequisites (editorial).}\begin{itemize}
\item Normalization commutes with field extensions up to universal homeomorphism. Stacks Project, Varieties, Lemma lemma-normalization-and-change-of-fields (tag 0C3N). Identifies the top horizontal arrow of the diagram and shows it is a universal homeomorphism.
\item Geometric branches as separable degree sums and connected components over separably closed fields. Stacks Project, Varieties, Lemma lemma-number-of-branches (tag 0C1S) and Lemma lemma-separably-closed-field-connected-components (tag 0363). The first gives the branch count along the normalization and the second converts the algebra fact into a connected-component count.
\end{itemize}\medskip
\ifdefined\hypertarget\hypertarget{sourcepointer-1}{}\fi
\noindent\textit{Source reference.} Original lemma 0C3N: sources/stacks-chow/tags.txt, source lines 15198--15198. Cross-chapter standard prerequisite cited by the author inside the counted span; it is used as a pointer and its own proof is not part of this item.\par\medskip
\ifdefined\hypertarget\hypertarget{sourcepointer-2}{}\fi
\noindent\textit{Source reference.} Original reference 04DF: sources/stacks-chow/tags.txt, source lines 5725--5725. Cross-chapter standard prerequisite cited by the author inside the counted span; it is used as a pointer and its own proof is not part of this item.\par\medskip
\ifdefined\hypertarget\hypertarget{sourcepointer-3}{}\fi
\noindent\textit{Source reference.} Original reference 01S4: sources/stacks-chow/tags.txt, source lines 2325--2325. Cross-chapter standard prerequisite cited by the author inside the counted span; it is used as a pointer and its own proof is not part of this item.\par\medskip
\ifdefined\hypertarget\hypertarget{sourcepointer-4}{}\fi
\noindent\textit{Source reference.} Original lemma 0C1S: sources/stacks-chow/tags.txt, source lines 15132--15132. Cross-chapter standard prerequisite cited by the author inside the counted span; it is used as a pointer and its own proof is not part of this item.\par\medskip
\ifdefined\hypertarget\hypertarget{sourcepointer-5}{}\fi
\noindent\textit{Source reference.} Original reference 01JT: sources/stacks-chow/tags.txt, source lines 2026--2026. Cross-chapter standard prerequisite cited by the author inside the counted span; it is used as a pointer and its own proof is not part of this item.\par\medskip
\ifdefined\hypertarget\hypertarget{sourcepointer-6}{}\fi
\noindent\textit{Source reference.} Original lemma 0363: sources/stacks-chow/tags.txt, source lines 4150--4150. Cross-chapter standard prerequisite cited by the author inside the counted span; it is used as a pointer and its own proof is not part of this item.\par\medskip
\begin{lemma}
\label{lemma-geometric-branches-and-change-of-fields}
Let $k$ be a field. Let $X$ be a locally algebraic $k$-scheme.
Let $K/k$ be an extension of fields. Let $y \in X_K$ be a
point with image $x$ in $X$. Then the number of
geometric branches of $X$ at $x$ is the number of geometric
branches of $X_K$ at $y$.
\end{lemma}

\begin{proof}
Write $Y = X_K$ and let $X^\nu$, resp.\ $Y^\nu$ be the normalization
of $X$, resp.\ $Y$. Consider the commutative diagram
$$
\xymatrix{
Y^\nu \ar[r] \ar[d] & X^\nu_K \ar[r] \ar[d]_{\nu_K} & X^\nu \ar[d]_\nu \\
Y \ar@{=}[r] & Y \ar[r] & X
}
$$
By Lemma \ref{lemma-normalization-and-change-of-fields} we see that
the left top horizontal arrow is a universal homeomorphism.
Hence it induces purely inseparable residue field extensions, see
Morphisms, Lemmas \ref{morphisms-lemma-universal-homeomorphism} and
\ref{morphisms-lemma-universally-injective}.
Thus the number of geometric branches of $Y$ at $y$ is
$\sum_{\nu_K(y') = y} [\kappa(y') : \kappa(y)]_s$
by Lemma \ref{lemma-number-of-branches}. Similarly
$\sum_{\nu(x') = x} [\kappa(x') : \kappa(x)]_s$ is the number of
geometric branches of $X$ at $x$. Using Schemes, Lemma
\ref{schemes-lemma-points-fibre-product}
our statement follows from the following algebra fact:
given a field extension $l/\kappa$ and
an algebraic field extension $m/\kappa$, then
$$
\sum\nolimits_{m \otimes_\kappa l \to m'} [m' : l']_s = [m : \kappa]_s
$$
where the sum is over the quotient fields of $m \otimes_\kappa l$.
One can prove this in an elementary way, or one can use
Lemma \ref{lemma-separably-closed-field-connected-components}
applied to
$$
\Spec(m \otimes_\kappa l) \times_{\Spec(l)} \Spec(\overline{l}) =
\Spec(m) \otimes_{\Spec(\kappa)} \Spec(\overline{l})
\longrightarrow
\Spec(m) \times_{\Spec(\kappa)} \Spec(\overline{\kappa})
$$
because one can interpret $[m : \kappa]_s$ as the number of
connected components of the right hand side and the sum
$\sum_{m \otimes_\kappa l \to m'} [m' : l']_s$
as the number of connected components of the left hand side.
\end{proof}

\end{document}
